\documentclass[../../../main.tex]{subfiles}

\begin{document}
% Exercise 1.4:2(d)

I would change the hypothesis of the exercise from saying that there are
infinitely many $k$ such that $n_k\neq 9$ to saying that there are
infinitely many $k$ such that $m_k\neq 9$.

The reason for this change is that the statement as written is false.
For example, consider
\[
m_0=1,\qquad m_k=9\qquad (k\geq 1).
\]
Then
\[
\sum_{i=0}^k m_i10^{-i}=2-10^{-k},
\]
so the supremum of these partial sums is $x=2$. Applying Rudin's
construction to $x$ gives
\[
n_0=2,\qquad n_k=0\qquad (k\geq 1).
\]
Thus $n_k\neq 9$ for infinitely many $k$, while $n_k\neq m_k$ for every
$k\geq 1$, contradicting the stated conclusion.

\begin{lemma}
\label{lemma:1.4.2.d.1}

    For every $k$,
    \[
        \sum_{i=0}^{k} m_i 10^{-i} + 10^{-k} > x.
    \]
\end{lemma}

\begin{proof}
	Suppose that
	\[
		\sum_{i=0}^{s} m_i 10^{-i} + 10^{-s} \leq x
	 \]
	for some $s$.
	Then, for any $t > s$,
	\[ 
		\sum_{i=0}^{s} m_i 10^{-i} + 10^{-s} - 10^{-t} < \sum_{i=0}^{s} m_i 10^{-i} + 10^{-s} \leq x.
	 \]
	Therefore, $\sum_{i=0}^{s} m_i 10^{-i} + 10^{-s} - 10^{-t}$ is not an upper bound of $\left\{ \sum_{i=0}^{k} m_i 10^{-i} \, \middle |\, k \ge 0 \right\}$.
	It follows that there exists $e \in \left\{ \sum_{i=0}^{k} m_i 10^{-i} \, \middle |\, k \ge 0 \right\}$ such that
	\[ 
		e > \sum_{i=0}^{s} m_i 10^{-i} + 10^{-s} - 10^{-t} = \sum_{i=0}^{s} m_i 10^{-i} + \sum_{i=s+1}^{t} 9 \cdot 10^{-i}.
	 \]
	This implies that $m_k = 9$ for every $k$ such that $s < k \leq t$.
		To see this, suppose that $m_j \neq 9$ for some $j \in \left(s, t\right]$,
	Then
	\[ 
		\sum_{i=0}^{j} m_i 10^{-i} + 10^{-j} \leq \sum_{i=0}^{s} m_i 10^{-i} + \sum_{i=s+1}^{t} 9 \cdot 10^{-i} < e.
	 \]
	But $\sum_{i=0}^{j} m_i 10^{-i} + 10^{-j} \geq e$ by part (c).

	Since $t > s$ can be chosen arbitrarily large, it follows that $m_k = 9$ for every $k > s$. Hence there are only finitely many $k$ such that $m_k \neq 9$, contradicting our hypothesis.
\end{proof}

Now we show by induction on $k$ that $n_k = m_k$ for all $k$.

\begin{proof}
	If $n_0 > m_0$, then
	\[
		n_0 \geq m_0 + 1 > x
	\]
	by Lemma \ref{lemma:1.4.2.d.1}.
	If $n_0 < m_0$, then $n_0$ is not the largest integer satisfying $n_0 \leq x$.
	Both cases contradict our choice of $n_0$. Therefore, $n_0 = m_0$.

	Now assume that $n_k = m_k$ for all $k < j$.
	If $n_j > m_j$, then
	\[
		\sum_{i=0}^{j} n_i 10^{-i}
		\geq
		\sum_{i=0}^{j} m_i 10^{-i} + 10^{-j}
		> x
	\]
	by Lemma \ref{lemma:1.4.2.d.1}.
	If $n_j < m_j$, then $n_j$ is not the largest integer satisfying
	\[
		\sum_{i=0}^{j} n_i 10^{-i} \leq x.
	\]
	Both cases contradict our choice of $n_j$. Therefore, $n_j = m_j$.

	Hence, by induction, $n_k = m_k$ for all $k$.
\end{proof}


\end{document}
